\SetPageTheme{story}
\PageHeader{Lectia 7}{Intelegerea algoritmului de oglindire}{L7 • P1}

\begin{missionbox}[O idee noua]
  Nu ne mai este suficient doar sa citim cifrele. Acum vrem sa construim un numar nou din ele.

  Cand verificam un palindrom, trebuie sa obtinem oglinditul numarului initial. Pentru asta nu ajunge sa afisam cifrele una cate una. Trebuie sa le asezam intr-o structura noua si coerenta.
\end{missionbox}

\begin{conceptbox}[Formula pentru oglindit]
  Daca \texttt{uc} este cifra extrasa, atunci:
  \[
    \texttt{ogl = ogl * 10 + uc}
  \]
  Mai intai mutam tot ce aveam in \texttt{ogl} o pozitie la stanga, apoi adaugam cifra noua.
\end{conceptbox}

\begin{guidebox}[Activare]
  Daca \texttt{ogl = 43} si \texttt{uc = 2}, dupa actualizare obtinem \hrulefill
\end{guidebox}

\newpage
\SetPageTheme{guided}
\PageHeader{Lectia 7}{Trace pe oglindire}{L7 • P2}

\begin{examplebox}[Exemplu pentru \texttt{n = 1234}]
  \begin{center}
  \begin{tabularx}{\linewidth}{|X|X|X|X|}
    \hline
    \textbf{Pas} & \textbf{uc} & \textbf{ogl} & \textbf{n dupa /10} \\
    \hline
    Start & - & 0 & 1234 \\
    \hline
    1 & 4 & 4 & 123 \\
    \hline
    2 & 3 & 43 & 12 \\
    \hline
    3 & 2 & 432 & 1 \\
    \hline
    4 & 1 & 4321 & 0 \\
    \hline
  \end{tabularx}
  \end{center}
\end{examplebox}

\begin{guidebox}[Exercitiu ghidat]
  Explica de ce formula adauga cifra noua la sfarsitul lui \texttt{ogl}.
  \AnswerSpace{3}
\end{guidebox}

\newpage
\SetPageTheme{guided}
\PageHeader{Lectia 7}{Semi-ghidat}{L7 • P3}

\begin{solobox}[Completeaza tabelul]
  Pentru \texttt{n = 508} completeaza:
  \begin{center}
  \begin{tabularx}{\linewidth}{|X|X|X|X|}
    \hline
    \textbf{Pas} & \textbf{uc} & \textbf{ogl} & \textbf{n dupa /10} \\
    \hline
    Start & - & 0 & 508 \\
    \hline
    1 & & & \\
    \hline
    2 & & & \\
    \hline
    3 & & & \\
    \hline
  \end{tabularx}
  \end{center}
\end{solobox}

\begin{solobox}[Independent scurt]
  Ce oglindit obtii pentru \texttt{n = 7002}? Scrie doar pasii esentiali.
  \AnswerSpace{4}
\end{solobox}

\newpage
\SetPageTheme{challenge}
\PageHeader{Lectia 7}{Independent}{L7 • P4}

\begin{solobox}[Fisa independenta]
  1. Scrie singur algoritmul care construieste oglinditul unui numar.

  2. Arata cum functioneaza pentru \texttt{n = 1221}.

  3. Explica de ce aceasta formula este mai puternica decat simpla afisare a cifrelor.
  \AnswerSpace{8}
\end{solobox}
